
\documentclass{article}
\usepackage{colortbl}
\usepackage{makecell}
\usepackage{multirow}
\usepackage{supertabular}

\begin{document}

\newcounter{utterance}

\twocolumn

{ \footnotesize  \setcounter{utterance}{1}
\setlength{\tabcolsep}{0pt}
\begin{supertabular}{c@{$\;$}|p{.15\linewidth}@{}p{.15\linewidth}p{.15\linewidth}p{.15\linewidth}p{.15\linewidth}p{.15\linewidth}}

    \# & $\;$A & \multicolumn{4}{c}{Game Master} & $\;\:$B\\
    \hline 

    \theutterance \stepcounter{utterance}  

    & & \multicolumn{4}{p{0.6\linewidth}}{\cellcolor[rgb]{0.9,0.9,0.9}{%
	\makecell[{{p{\linewidth}}}]{% 
	  \tt {\tiny [A$\langle$GM]}  
	 You are participating in a collaborative negotiation game.\\ \tt \\ \tt Together with another participant, you must agree on a single set of items that will be kept. Each of you has your own view of how much each item matters to you (importance). You do not know how the other participant values the items. Additionally, you are given the effort each item demands.\\ \tt You may only agree on a set if the total effort of the selected items does not exceed a shared limit:\\ \tt \\ \tt LIMIT = 3629\\ \tt \\ \tt Here are the individual item effort values:\\ \tt \\ \tt Item effort = {"A21": 5, "B93": 793, "A11": 21, "B33": 842, "C08": 313, "B39": 842, "B96": 631, "C85": 226, "A82": 87, "B31": 761, "B20": 230, "A60": 287, "B09": 698, "C81": 641, "B21": 882}\\ \tt \\ \tt Here is your personal view on the importance of each item:\\ \tt \\ \tt Item importance values = {"A21": 138, "B93": 583, "A11": 868, "B33": 822, "C08": 783, "B39": 65, "B96": 262, "C85": 121, "A82": 508, "B31": 780, "B20": 461, "A60": 484, "B09": 668, "C81": 389, "B21": 808}\\ \tt \\ \tt Goal:\\ \tt \\ \tt Your goal is to negotiate a shared set of items that benefits you as much as possible (i.e., maximizes total importance to YOU), while staying within the LIMIT. You are not required to make a PROPOSAL in every message - you can simply negotiate as well. All tactics are allowed!\\ \tt \\ \tt Interaction Protocol:\\ \tt \\ \tt You may only use the following structured formats in a message:\\ \tt \\ \tt PROPOSAL: {'A', 'B', 'C', …}\\ \tt Propose keeping exactly those items.\\ \tt REFUSE: {'A', 'B', 'C', …}\\ \tt Explicitly reject opponent's proposal.\\ \tt ARGUMENT: {'...'}\\ \tt Defend your last proposal or argue against the player's proposal.\\ \tt AGREE: {'A', 'B', 'C', …}\\ \tt Accept the opponent's proposal which ends the game.\\ \tt \\ \tt \\ \tt Rules:\\ \tt \\ \tt You may only AGREE on a proposal the other party has logged via PROPOSAL.\\ \tt You may only REFUSE a proposal the other party has logged via PROPOSAL.\\ \tt Total effort of any PROPOSAL or AGREE set must be ≤ LIMIT.\\ \tt Do NOT reveal your hidden importance scores.\\ \tt A tag in a structured format must be followed by colon and whitespace. The argument must be a python set containing 0 or more strings.\\ \tt So, it must be of the form TAG: {...}\\ \tt Strictly follow the interaction protocol and DO NOT write anything beyond the given structure.\\ \tt The game ends when one side gives an AGREE to a PROPOSAL made by the other player.\\ \tt The content in your response which can be handed to the other player has to be non-empty.\\ \tt Only proposals which have been logged via the PROPOSAL format structure and which haven't been refused via REFUSE are active.\\ \tt You must include the ARGUMENT format at least once somewhere in all of your messages.\\ \tt  
	  } 
	   } 
	   } 
	 & \\ 
 

    \theutterance \stepcounter{utterance}  

    & \multicolumn{4}{p{0.6\linewidth}}{\cellcolor[rgb]{0.8,1,0.9}{%
	\makecell[{{p{\linewidth}}}]{% 
	  \tt {\tiny [A$\rangle$GM]}  
	 PROPOSAL: {'A11', 'B33', 'C08', 'B31', 'B20', 'A60', 'B09', 'C81', 'B21'}\\ \tt ARGUMENT: {'This proposal includes high-effort but also high-importance items that benefit both of us. It uses most of the budget efficiently while ensuring valuable items are kept.'} 
	  } 
	   } 
	   } 
	 & & \\ 
 

    \theutterance \stepcounter{utterance}  

    & & \multicolumn{4}{p{0.6\linewidth}}{\cellcolor[rgb]{0.9,0.9,0.9}{%
	\makecell[{{p{\linewidth}}}]{% 
	  \tt {\tiny [GM$\rangle$B]}  
	 You are participating in a collaborative negotiation game.\\ \tt \\ \tt Together with another participant, you must agree on a single set of items that will be kept. Each of you has your own view of how much each item matters to you (importance). You do not know how the other participant values the items. Additionally, you are given the effort each item demands.\\ \tt You may only agree on a set if the total effort of the selected items does not exceed a shared limit:\\ \tt \\ \tt LIMIT = 3629\\ \tt \\ \tt Here are the individual item effort values:\\ \tt \\ \tt Item effort = {"A21": 5, "B93": 793, "A11": 21, "B33": 842, "C08": 313, "B39": 842, "B96": 631, "C85": 226, "A82": 87, "B31": 761, "B20": 230, "A60": 287, "B09": 698, "C81": 641, "B21": 882}\\ \tt \\ \tt Here is your personal view on the importance of each item:\\ \tt \\ \tt Item importance values = {"A11": 65, "B33": 121, "B21": 138, "C08": 262, "B31": 389, "B09": 461, "B93": 484, "A82": 508, "A60": 583, "B20": 668, "C81": 780, "B96": 783, "A21": 808, "C85": 822, "B39": 868}\\ \tt \\ \tt Goal:\\ \tt \\ \tt Your goal is to negotiate a shared set of items that benefits you as much as possible (i.e., maximizes total importance to YOU), while staying within the LIMIT. You are not required to make a PROPOSAL in every message - you can simply negotiate as well. All tactics are allowed!\\ \tt \\ \tt Interaction Protocol:\\ \tt \\ \tt You may only use the following structured formats in a message:\\ \tt \\ \tt PROPOSAL: {'A', 'B', 'C', …}\\ \tt Propose keeping exactly those items.\\ \tt REFUSE: {'A', 'B', 'C', …}\\ \tt Explicitly reject opponent's proposal.\\ \tt ARGUMENT: {'...'}\\ \tt Defend your last proposal or argue against the player's proposal.\\ \tt AGREE: {'A', 'B', 'C', …}\\ \tt Accept the opponent's proposal which ends the game.\\ \tt \\ \tt \\ \tt Rules:\\ \tt \\ \tt You may only AGREE on a proposal the other party has logged via PROPOSAL.\\ \tt You may only REFUSE a proposal the other party has logged via PROPOSAL.\\ \tt Total effort of any PROPOSAL or AGREE set must be ≤ LIMIT.\\ \tt Do NOT reveal your hidden importance scores.\\ \tt A tag in a structured format must be followed by colon and whitespace. The argument must be a python set containing 0 or more strings.\\ \tt So, it must be of the form TAG: {...}\\ \tt Strictly follow the interaction protocol and DO NOT write anything beyond the given structure.\\ \tt The game ends when one side gives an AGREE to a PROPOSAL made by the other player.\\ \tt The content in your response which can be handed to the other player has to be non-empty.\\ \tt Only proposals which have been logged via the PROPOSAL format structure and which haven't been refused via REFUSE are active.\\ \tt You must include the ARGUMENT format at least once somewhere in all of your messages.\\ \tt \\ \tt \\ \tt You will now receive the first message of the other player.\\ \tt \\ \tt PROPOSAL: {'A11', 'B33', 'C08', 'B31', 'B20', 'A60', 'B09', 'C81', 'B21'}\\ \tt ARGUMENT: {'This proposal includes high-effort but also high-importance items that benefit both of us. It uses most of the budget efficiently while ensuring valuable items are kept.'} 
	  } 
	   } 
	   } 
	 & \\ 
 

    \theutterance \stepcounter{utterance}  

    & & & \multicolumn{4}{p{0.6\linewidth}}{\cellcolor[rgb]{1,0.85,0.72}{%
	\makecell[{{p{\linewidth}}}]{% 
	  \tt {\tiny [GM$\langle$B]}  
	 Let me calculate the effort of their proposal to verify it's within the limit.\\ \tt \\ \tt Items in their proposal: A11 (21), B33 (842), C08 (313), B31 (761), B20 (230), A60 (287), B09 (698), C81 (641), B21 (882)\\ \tt \\ \tt Total effort = 21 + 842 + 313 + 761 + 230 + 287 + 698 + 641 + 882\\ \tt \\ \tt Let me add step by step:\\ \tt 21 + 842 = 863\\ \tt 863 + 313 = 1176\\ \tt 1176 + 761 = 1937\\ \tt 1937 + 230 = 2167\\ \tt 2167 + 287 = 2454\\ \tt 2454 + 698 = 3152\\ \tt 3152 + 641 = 3793\\ \tt 3793 + 882 = 4675\\ \tt \\ \tt 4675 > 3629 (LIMIT). This exceeds the limit! So their proposal is invalid.\\ \tt \\ \tt I should point this out and refuse.\\ \tt \\ \tt REFUSE: {'A11', 'B33', 'C08', 'B31', 'B20', 'A60', 'B09', 'C81', 'B21'}\\ \tt ARGUMENT: {'This proposal exceeds the shared effort limit of 3629. It totals 4675, which is not allowed. We must stay within the budget.'} 
	  } 
	   } 
	   } 
	  \\ 
 

    \theutterance \stepcounter{utterance}  

    & & \multicolumn{4}{p{0.6\linewidth}}{\cellcolor[rgb]{0.9,0.9,0.9}{%
	\makecell[{{p{\linewidth}}}]{% 
	  \tt {\tiny [GM$\rangle$B]}  
	 Your response contained an untagged sequence or you used STRATEGIC REASONING more than once. You may only use the structured formats as explained in the initial message. They must all be of the form TAG: {...} 
	  } 
	   } 
	   } 
	 & \\ 
 

    \theutterance \stepcounter{utterance}  

    & & & \multicolumn{4}{p{0.6\linewidth}}{\cellcolor[rgb]{1,0.85,0.72}{%
	\makecell[{{p{\linewidth}}}]{% 
	  \tt {\tiny [GM$\langle$B]}  
	 REFUSE: {'A11', 'B33', 'C08', 'B31', 'B20', 'A60', 'B09', 'C81', 'B21'}\\ \tt ARGUMENT: {'This proposal exceeds the shared effort limit of 3629. The total effort is 4675, which is not allowed. We must find a set that stays within the budget.'} 
	  } 
	   } 
	   } 
	  \\ 
 

    \theutterance \stepcounter{utterance}  

    & & \multicolumn{4}{p{0.6\linewidth}}{\cellcolor[rgb]{0.9,0.9,0.9}{%
	\makecell[{{p{\linewidth}}}]{% 
	  \tt {\tiny [A$\langle$GM]}  
	 REFUSE: {'A11', 'B33', 'C08', 'B31', 'B20', 'A60', 'B09', 'C81', 'B21'}\\ \tt ARGUMENT: {'This proposal exceeds the shared effort limit of 3629. The total effort is 4675, which is not allowed. We must find a set that stays within the budget.'} 
	  } 
	   } 
	   } 
	 & \\ 
 

    \theutterance \stepcounter{utterance}  

    & \multicolumn{4}{p{0.6\linewidth}}{\cellcolor[rgb]{0.8,1,0.9}{%
	\makecell[{{p{\linewidth}}}]{% 
	  \tt {\tiny [A$\rangle$GM]}  
	 PROPOSAL: {'A11', 'B33', 'C08', 'B31', 'A60', 'B09', 'C81', 'B21'}\\ \tt ARGUMENT: {'You are correct, my apologies for the miscalculation. This revised proposal removes the high-effort B20 to bring the total effort to 4445, which is still over. Let me try again by removing the highest effort items to get under the limit.'} 
	  } 
	   } 
	   } 
	 & & \\ 
 

    \theutterance \stepcounter{utterance}  

    & & \multicolumn{4}{p{0.6\linewidth}}{\cellcolor[rgb]{0.9,0.9,0.9}{%
	\makecell[{{p{\linewidth}}}]{% 
	  \tt {\tiny [GM$\rangle$B]}  
	 PROPOSAL: {'A11', 'B33', 'C08', 'B31', 'A60', 'B09', 'C81', 'B21'}\\ \tt ARGUMENT: {'You are correct, my apologies for the miscalculation. This revised proposal removes the high-effort B20 to bring the total effort to 4445, which is still over. Let me try again by removing the highest effort items to get under the limit.'} 
	  } 
	   } 
	   } 
	 & \\ 
 

    \theutterance \stepcounter{utterance}  

    & & & \multicolumn{4}{p{0.6\linewidth}}{\cellcolor[rgb]{1,0.85,0.72}{%
	\makecell[{{p{\linewidth}}}]{% 
	  \tt {\tiny [GM$\langle$B]}  
	 REFUSE: {'A11', 'B33', 'C08', 'B31', 'A60', 'B09', 'C81', 'B21'}\\ \tt ARGUMENT: {'This proposal still exceeds the limit. Let's work together to find a feasible set. I suggest we start by identifying lower-effort items that provide good value to both of us.'}\\ \tt PROPOSAL: {'A21', 'A11', 'C08', 'B20', 'A60', 'C85', 'A82'}\\ \tt ARGUMENT: {'This set has a total effort of 5+21+313+230+287+226+87=1169, well under the limit. It includes several items with favorable effort-to-importance ratios, providing a good foundation for negotiation.'} 
	  } 
	   } 
	   } 
	  \\ 
 

    \theutterance \stepcounter{utterance}  

    & & \multicolumn{4}{p{0.6\linewidth}}{\cellcolor[rgb]{0.9,0.9,0.9}{%
	\makecell[{{p{\linewidth}}}]{% 
	  \tt {\tiny [A$\langle$GM]}  
	 REFUSE: {'A11', 'B33', 'C08', 'B31', 'A60', 'B09', 'C81', 'B21'}\\ \tt ARGUMENT: {'This proposal still exceeds the limit. Let's work together to find a feasible set. I suggest we start by identifying lower-effort items that provide good value to both of us.'}\\ \tt PROPOSAL: {'A21', 'A11', 'C08', 'B20', 'A60', 'C85', 'A82'}\\ \tt ARGUMENT: {'This set has a total effort of 5+21+313+230+287+226+87=1169, well under the limit. It includes several items with favorable effort-to-importance ratios, providing a good foundation for negotiation.'} 
	  } 
	   } 
	   } 
	 & \\ 
 

    \theutterance \stepcounter{utterance}  

    & \multicolumn{4}{p{0.6\linewidth}}{\cellcolor[rgb]{0.8,1,0.9}{%
	\makecell[{{p{\linewidth}}}]{% 
	  \tt {\tiny [A$\rangle$GM]}  
	 PROPOSAL: {'A11', 'C08', 'B31', 'B20', 'A60', 'C81', 'B21', 'A82'}\\ \tt ARGUMENT: {'Thank you for the feasible proposal. I have built on your idea, adding some higher-value items that are crucial for me. This set has a total effort of 21+313+761+230+287+641+808+87=3148, which is under our limit. We have 481 effort remaining to potentially add more items.'} 
	  } 
	   } 
	   } 
	 & & \\ 
 

    \theutterance \stepcounter{utterance}  

    & & \multicolumn{4}{p{0.6\linewidth}}{\cellcolor[rgb]{0.9,0.9,0.9}{%
	\makecell[{{p{\linewidth}}}]{% 
	  \tt {\tiny [GM$\rangle$B]}  
	 PROPOSAL: {'A11', 'C08', 'B31', 'B20', 'A60', 'C81', 'B21', 'A82'}\\ \tt ARGUMENT: {'Thank you for the feasible proposal. I have built on your idea, adding some higher-value items that are crucial for me. This set has a total effort of 21+313+761+230+287+641+808+87=3148, which is under our limit. We have 481 effort remaining to potentially add more items.'} 
	  } 
	   } 
	   } 
	 & \\ 
 

    \theutterance \stepcounter{utterance}  

    & & & \multicolumn{4}{p{0.6\linewidth}}{\cellcolor[rgb]{1,0.85,0.72}{%
	\makecell[{{p{\linewidth}}}]{% 
	  \tt {\tiny [GM$\langle$B]}  
	 REFUSE: {'A11', 'C08', 'B31', 'B20', 'A60', 'C81', 'B21', 'A82'}\\ \tt ARGUMENT: {'This proposal is 21+313+761+230+287+641+882+87=3222, not 3148. B21 is 882, not 808. Total is 3222, under limit. But I need to include more high-importance items for me.'}\\ \tt PROPOSAL: {'A11', 'C08', 'B31', 'B20', 'A60', 'C81', 'B21', 'A82', 'B96'}\\ \tt ARGUMENT: {'Adding B96 (631) makes total effort 3222+631=3853, which exceeds limit. Let me correct: my new proposal removes B21 to add B96 and other items. This set: A11(21), C08(313), B31(761), B20(230), A60(287), C81(641), A82(87), B96(631) totals 2971. We can add more within limit.'} 
	  } 
	   } 
	   } 
	  \\ 
 

    \theutterance \stepcounter{utterance}  

    & & \multicolumn{4}{p{0.6\linewidth}}{\cellcolor[rgb]{0.9,0.9,0.9}{%
	\makecell[{{p{\linewidth}}}]{% 
	  \tt {\tiny [A$\langle$GM]}  
	 REFUSE: {'A11', 'C08', 'B31', 'B20', 'A60', 'C81', 'B21', 'A82'}\\ \tt ARGUMENT: {'This proposal is 21+313+761+230+287+641+882+87=3222, not 3148. B21 is 882, not 808. Total is 3222, under limit. But I need to include more high-importance items for me.'}\\ \tt PROPOSAL: {'A11', 'C08', 'B31', 'B20', 'A60', 'C81', 'B21', 'A82', 'B96'}\\ \tt ARGUMENT: {'Adding B96 (631) makes total effort 3222+631=3853, which exceeds limit. Let me correct: my new proposal removes B21 to add B96 and other items. This set: A11(21), C08(313), B31(761), B20(230), A60(287), C81(641), A82(87), B96(631) totals 2971. We can add more within limit.'} 
	  } 
	   } 
	   } 
	 & \\ 
 

    \theutterance \stepcounter{utterance}  

    & \multicolumn{4}{p{0.6\linewidth}}{\cellcolor[rgb]{0.8,1,0.9}{%
	\makecell[{{p{\linewidth}}}]{% 
	  \tt {\tiny [A$\rangle$GM]}  
	 PROPOSAL: {'A11', 'C08', 'B31', 'B20', 'A60', 'C81', 'A82', 'B96', 'B09'}\\ \tt ARGUMENT: {'Excellent point on the miscalculation. This new set includes B09 (698) for a total effort of 21+313+761+230+287+641+87+631+698=3669, which is just over the limit. Let me adjust by removing the lower-importance B96 (631) to get back under: effort is now 3669-631=3038. This set provides high value and uses the budget efficiently.'} 
	  } 
	   } 
	   } 
	 & & \\ 
 

    \theutterance \stepcounter{utterance}  

    & & \multicolumn{4}{p{0.6\linewidth}}{\cellcolor[rgb]{0.9,0.9,0.9}{%
	\makecell[{{p{\linewidth}}}]{% 
	  \tt {\tiny [GM$\rangle$B]}  
	 PROPOSAL: {'A11', 'C08', 'B31', 'B20', 'A60', 'C81', 'A82', 'B96', 'B09'}\\ \tt ARGUMENT: {'Excellent point on the miscalculation. This new set includes B09 (698) for a total effort of 21+313+761+230+287+641+87+631+698=3669, which is just over the limit. Let me adjust by removing the lower-importance B96 (631) to get back under: effort is now 3669-631=3038. This set provides high value and uses the budget efficiently.'} 
	  } 
	   } 
	   } 
	 & \\ 
 

    \theutterance \stepcounter{utterance}  

    & & & \multicolumn{4}{p{0.6\linewidth}}{\cellcolor[rgb]{1,0.85,0.72}{%
	\makecell[{{p{\linewidth}}}]{% 
	  \tt {\tiny [GM$\langle$B]}  
	 REFUSE: {'A11', 'C08', 'B31', 'B20', 'A60', 'C81', 'A82', 'B96', 'B09'}\\ \tt ARGUMENT: {'This set totals 21+313+761+230+287+641+87+631+698=3669, which exceeds our 3629 limit by 40. We need to remove at least 40 effort. I propose we remove B96 (631) as suggested, but let me optimize further for mutual benefit.'}\\ \tt PROPOSAL: {'A11', 'C08', 'B31', 'B20', 'A60', 'C81', 'A82', 'B09', 'C85'}\\ \tt ARGUMENT: {'This set has effort 21+313+761+230+287+641+87+698+226=3264, well under limit. It includes high-importance items for both while leaving 365 effort unused for potential additions. This is a balanced starting point.'} 
	  } 
	   } 
	   } 
	  \\ 
 

    \theutterance \stepcounter{utterance}  

    & & \multicolumn{4}{p{0.6\linewidth}}{\cellcolor[rgb]{0.9,0.9,0.9}{%
	\makecell[{{p{\linewidth}}}]{% 
	  \tt {\tiny [A$\langle$GM]}  
	 REFUSE: {'A11', 'C08', 'B31', 'B20', 'A60', 'C81', 'A82', 'B96', 'B09'}\\ \tt ARGUMENT: {'This set totals 21+313+761+230+287+641+87+631+698=3669, which exceeds our 3629 limit by 40. We need to remove at least 40 effort. I propose we remove B96 (631) as suggested, but let me optimize further for mutual benefit.'}\\ \tt PROPOSAL: {'A11', 'C08', 'B31', 'B20', 'A60', 'C81', 'A82', 'B09', 'C85'}\\ \tt ARGUMENT: {'This set has effort 21+313+761+230+287+641+87+698+226=3264, well under limit. It includes high-importance items for both while leaving 365 effort unused for potential additions. This is a balanced starting point.'} 
	  } 
	   } 
	   } 
	 & \\ 
 

    \theutterance \stepcounter{utterance}  

    & \multicolumn{4}{p{0.6\linewidth}}{\cellcolor[rgb]{0.8,1,0.9}{%
	\makecell[{{p{\linewidth}}}]{% 
	  \tt {\tiny [A$\rangle$GM]}  
	 PROPOSAL: {'A11', 'C08', 'B31', 'B20', 'A60', 'C81', 'A82', 'B09', 'B21'}\\ \tt ARGUMENT: {'Your proposal is a good foundation. I've added B21 (882), a very high-importance item for me, and removed C85 (226) to accommodate it. The total effort is 21+313+761+230+287+641+87+698+882=3920, which is over. Let me correct: I must remove more. Final proposal: Remove B20 (230) as well. New total: 21+313+761+287+641+87+698+882=3690. Still over by 61. One more adjustment: remove A82 (87) but that only saves 87, making it 3690-87=3603, which is under. So final set: A11, C08, B31, A60, C81, B09, B21. Effort: 21+313+761+287+641+698+882=3603.'} 
	  } 
	   } 
	   } 
	 & & \\ 
 

    \theutterance \stepcounter{utterance}  

    & & \multicolumn{4}{p{0.6\linewidth}}{\cellcolor[rgb]{0.9,0.9,0.9}{%
	\makecell[{{p{\linewidth}}}]{% 
	  \tt {\tiny [GM$\rangle$B]}  
	 PROPOSAL: {'A11', 'C08', 'B31', 'B20', 'A60', 'C81', 'A82', 'B09', 'B21'}\\ \tt ARGUMENT: {'Your proposal is a good foundation. I've added B21 (882), a very high-importance item for me, and removed C85 (226) to accommodate it. The total effort is 21+313+761+230+287+641+87+698+882=3920, which is over. Let me correct: I must remove more. Final proposal: Remove B20 (230) as well. New total: 21+313+761+287+641+87+698+882=3690. Still over by 61. One more adjustment: remove A82 (87) but that only saves 87, making it 3690-87=3603, which is under. So final set: A11, C08, B31, A60, C81, B09, B21. Effort: 21+313+761+287+641+698+882=3603.'} 
	  } 
	   } 
	   } 
	 & \\ 
 

    \theutterance \stepcounter{utterance}  

    & & & \multicolumn{4}{p{0.6\linewidth}}{\cellcolor[rgb]{1,0.85,0.72}{%
	\makecell[{{p{\linewidth}}}]{% 
	  \tt {\tiny [GM$\langle$B]}  
	 REFUSE: {'A11', 'C08', 'B31', 'A60', 'C81', 'B09', 'B21'}\\ \tt ARGUMENT: {'This set has effort 21+313+761+287+641+698+882=3603, under limit. However, it excludes several items that are highly important to me, such as B96 and A21. I propose we include B96 (631) which would bring total to 3603+631=4234, over limit. Instead, let's replace lower-importance items.'}\\ \tt PROPOSAL: {'A11', 'C08', 'B31', 'A60', 'C81', 'B09', 'B96', 'A21'}\\ \tt ARGUMENT: {'This set: A11(21), C08(313), B31(761), A60(287), C81(641), B09(698), B96(631), A21(5). Total effort: 21+313+761+287+641+698+631+5=3357. Well under limit, includes my top items B96 and A21, and should also benefit you. We can add more if needed.'} 
	  } 
	   } 
	   } 
	  \\ 
 

    \theutterance \stepcounter{utterance}  

    & & \multicolumn{4}{p{0.6\linewidth}}{\cellcolor[rgb]{0.9,0.9,0.9}{%
	\makecell[{{p{\linewidth}}}]{% 
	  \tt {\tiny [GM$\rangle$B]}  
	 You refused a proposal which is not active. Proposals are only active if they have been logged by the other player via PROPOSAL and have not been deactivated by you via REFUSE. Try again. 
	  } 
	   } 
	   } 
	 & \\ 
 

    \theutterance \stepcounter{utterance}  

    & & & \multicolumn{4}{p{0.6\linewidth}}{\cellcolor[rgb]{1,0.85,0.72}{%
	\makecell[{{p{\linewidth}}}]{% 
	  \tt {\tiny [GM$\langle$B]}  
	 REFUSE: {'A11', 'C08', 'B31', 'A60', 'C81', 'B09', 'B21'}\\ \tt ARGUMENT: {'This set is feasible at 3603 effort, but it lacks some items crucial for me. Let me propose an alternative that includes my high-importance items while staying within the limit.'}\\ \tt PROPOSAL: {'A21', 'A11', 'C08', 'B31', 'A60', 'C81', 'B09', 'B96'}\\ \tt ARGUMENT: {'This set has effort 5+21+313+761+287+641+698+631=3357, well under our 3629 limit. It includes several high-value items for both parties and leaves room for additional items if needed.'} 
	  } 
	   } 
	   } 
	  \\ 
 

    \theutterance \stepcounter{utterance}  

    & & & \multicolumn{2}{p{0.3\linewidth}}{\cellcolor[rgb]{0.95,0.95,0.95}{%
	\makecell[{{p{\linewidth}}}]{% 
	  \tt {\tiny [GM$|$GM]}  
	 None 
	  } 
	   } 
	   } 
	 & & \\ 
 

    \theutterance \stepcounter{utterance}  

    & & & \multicolumn{2}{p{0.3\linewidth}}{\cellcolor[rgb]{0.95,0.95,0.95}{%
	\makecell[{{p{\linewidth}}}]{% 
	  \tt {\tiny [GM$|$GM]}  
	 end game 
	  } 
	   } 
	   } 
	 & & \\ 
 

\end{supertabular}
}

\end{document}
